\section{文化与思想：在秩序失灵中重新问“什么是对的”}

春秋晚期，礼不再只是稳定身份的程序，也成为人们批评现实的语言。孔子及其后的儒者不是在一个完整的周礼世界中说话，而是在旧秩序还能被引用、却越来越难被实行的时代里寻找答案。

鲁国保存的年表、诸侯国之间不断重演的会盟和背约、卿族侵蚀国君权力的过程，都让“礼”显出两层意义：它既是现实中仍被使用的仪式，也是一种越来越像理想尺度的语言。人们不必先拥有一个完整哲学体系，才会开始问“什么是合乎礼的”；恰恰是在旧制度无法自动解决冲突时，这个问题才变得尖锐。

孔子周游列国的经历常被讲成一位圣人的孤独旅行。放回当时，它也是一个失去稳定政治位置的士人，试图在鲁、卫、陈、蔡等国之间寻找可以实行自己主张的空间。各国君主愿意听他说话，不等于愿意承担改变权力分配的代价。思想的挫折因此不是人物性格的悲剧，而是春秋末年政治结构的结果。

\begin{textwindow}[“礼乐征伐自诸侯出”】【论语】]
《论语》说：“天下有道，则礼乐征伐自天子出；天下无道，则礼乐征伐自诸侯出。”它不是一份中立新闻，而是后人对政治权力下移的概括。放在春秋章节里，能帮助读者理解为什么周王室的名分和诸侯的实力会长期错位。
\end{textwindow}

\subsection{《春秋》为什么既短又重}

《春秋》的文字短得近乎冷淡：会盟、出兵、弑君、灾异，经常只用十几个字。它之所以后来变得沉重，不是因为每个字天然包含秘密，而是因为《左传》《公羊传》《谷梁传》以及后来的史家不断追问：为什么这样记，而不那样记？谁被写进年表，谁被省略，哪种称呼带着褒贬？

这也是历史书写本身的一条线。春秋的人正在经历秩序的松动；后来的读者又用这部年表寻找秩序为何松动、应当怎样恢复。孔子、左传作者、汉代经师和现代史家读到的，并不是完全相同的“春秋”。
